\documentclass[11pt]{article}

\usepackage[T1]{fontenc}
\usepackage[utf8]{inputenc}
\usepackage{lmodern}
\usepackage{microtype}
\usepackage[a4paper,margin=0.85in]{geometry}
\usepackage{amsmath,amssymb}
\usepackage{booktabs}
\usepackage{tabularx}
\usepackage{longtable}
\usepackage{array}
\usepackage{graphicx}
\usepackage{caption}
\usepackage{subcaption}
\usepackage{float}
\usepackage{enumitem}
\usepackage{xcolor}
\usepackage{url}
\usepackage{hyperref}
\usepackage{fancyhdr}

\graphicspath{{../rlhf_runs/qwen25_05b_helpsteer3_trl_a100_full/full/eval_ppo_guarded_r768_b64_kl10_rm_ep2_u188_full2017/plots/}}
\urlstyle{same}
\hypersetup{
  colorlinks=true,
  linkcolor=blue!45!black,
  citecolor=blue!45!black,
  urlcolor=blue!55!black,
  pdftitle={Technical Companion: Post-Training Qwen2.5-0.5B with RLHF},
  pdfauthor={Dheeraj Dhillon}
}

\pagestyle{fancy}
\fancyhf{}
\lhead{RLHF Technical Companion}
\rhead{\thepage}
\setlength{\headheight}{14pt}
\setlist{nosep,leftmargin=1.4em}
\setlength{\parindent}{0pt}
\setlength{\parskip}{0.55em}
\captionsetup{font=small,labelfont=bf}

\newcommand{\model}{\textsc{Qwen2.5-0.5B-Instruct}}
\newcommand{\dataset}{\textsc{HelpSteer3-Preference}}
\newcommand{\trl}{\textsc{TRL}}
\newcommand{\projecturl}{\url{https://github.com/djdhillxn/rlhf}}
\newcommand{\explorerurl}{\url{https://djdhillxn.github.io/projects/rlhf-comparison/}}
\newcolumntype{Y}{>{\raggedright\arraybackslash}X}
\newcolumntype{R}{>{\raggedleft\arraybackslash}p{1.55cm}}

\title{
  \vspace{-1.5cm}
  \textbf{Technical Companion and Experiment Record}\\
  \large Post-Training Qwen2.5-0.5B with SFT, Reward Modeling, Guarded PPO, and Archived Extensions
}
\author{Dheeraj Dhillon}
\date{July 2026}

\begin{document}

\maketitle

\begin{abstract}
This companion preserves the engineering, experimental, and artifact-level
record behind the shorter two-column project report. It is intentionally more
operational: it records the custom-to-\trl{} migration, sequence-length studies,
reward-model audits, PPO failure modes, checkpoint and Colab contracts, GPU
memory redesign, full-suite evaluation, qualitative curation, and the
implemented-but-unexecuted DPO extension. The concise project report remains
the primary academic account; this document is the durable experiment archive.

The final pipeline trains \model{} on NVIDIA's heterogeneous \dataset{}. A
4,096-token SFT and reward-model budget retains almost all long-form examples.
One epoch of response-only LoRA SFT reaches 72.02\% validation mean token
accuracy, and an SFT-initialized two-epoch Bradley--Terry reward model reaches
65.62\% pairwise accuracy on 1,917 held-out pairs. The selected policy is then
trained for 12,032 on-policy rollouts and 188 PPO updates with 768-token
responses, domain-balanced sampling, quantile-calibrated reward bounds,
EOS-aware terminal scoring, smooth repeated-4-gram shaping, KL control, and an
exact-resume contract.

The completed 2,017-prompt suite is deliberately mixed. The learned reward
model prefers guarded PPO over Base in 817 comparisons (40.51\%) and over SFT
in 974 (48.29\%). PPO nevertheless completes without empty-output or numerical
collapse, has a 16.31\% cap-hit rate, and crosses a 25\% repeated-4-gram
threshold in 16.16\% of responses. An earlier 100-update TRL checkpoint scored
better under the same learned proxy but was substantially longer, more capped,
and more repetitive; it is retained here as evidence of reward-model
exploitation rather than presented as the final result. The main contribution
is therefore an inspectable RLHF system and an honest account of the tradeoff
between proxy optimization, behavioral control, and qualitative response
quality.
\end{abstract}

\tableofcontents
\newpage

\section{Introduction}

Large language models are trained primarily by next-token prediction, while
users ultimately care about properties such as helpfulness, relevance,
correctness, clarity, and safety. RLHF addresses this mismatch by learning a
reward signal from preference comparisons and then optimizing a language model
against that signal \cite{christiano2017,stiennon2020,ouyang2022}. In its
canonical form, the procedure has three stages: supervised fine-tuning on
preferred demonstrations, training a reward model on chosen/rejected pairs,
and policy optimization with a KL constraint to a reference policy.

% This project asks a deliberately modest but technically demanding question:
% can the complete procedure be made measurable and inspectable on a
% half-billion-parameter instruction model? The small policy makes iteration
% possible on rented notebook hardware, but it also makes the experiment a
% stress test of data quality, reward reliability, implementation detail, and
% evaluation discipline. The base model, \model{}, is already instruction-tuned
% and belongs to the Qwen2.5 family \cite{qwen25}. The work is therefore
% preference post-training rather than instruction tuning from an unaligned base.

This project deliberately pairs a modest compute budget with a nontrivial
preference distribution. It asks whether the complete procedure can be made
measurable and inspectable on a half-billion-parameter instruction model while
retaining HelpSteer3's general, code, STEM, multilingual, and long-form
examples. The small policy makes iteration possible on rented notebook
hardware, but it also makes the experiment a stress test of data quality,
reward reliability, implementation detail, and evaluation discipline. The
base model, \model{}, is already instruction-tuned and belongs to the Qwen2.5
family \cite{qwen25}. The work is therefore preference post-training rather
than instruction tuning from an unaligned base.

The project also connects two settings that are often discussed separately.
Earlier work in the same codebase implemented Trust Region Policy Optimization
(TRPO), Natural Policy Gradient, and PPO for MuJoCo and Atari. TRPO constrains
policy updates through a trust region \cite{schulman2015}; PPO approximates the
same conservative-update principle with a clipped surrogate objective
\cite{schulman2017}. In language-model RLHF, a second trust-region-like
constraint appears through the KL penalty to a frozen SFT reference. The
optimization machinery changes scale, but the central concern remains the
same: improve expected reward without allowing destructive policy drift.

\subsection{Contributions}

The project makes six practical contributions:

\begin{enumerate}
  \item It implements the full SFT--reward-model--PPO pipeline first with
  custom training code and then with \trl{}, preserving the custom lineage as
  a reproducible record of the assumptions and failures that motivated the
  final design.
  \item It develops response-safe, long-context preprocessing for
  \dataset{}, including assistant-only loss masks, complete EOS-terminated
  completions, distinct PAD/EOS tokens, prompt-first truncation, and
  prompt-identity tracking across training and evaluation.
  \item It initializes reward and value learning from the supervised policy
  and adopts high-impact PPO details such as zero dropout, matched behavior
  probabilities, fixed-length rollout masks, explicit missing-EOS treatment,
  and separate policy/reference semantics.
  \item It adds guarded long-form PPO: balanced domain sampling,
  quantile-calibrated reward clipping, lower-bound scoring for missing EOS,
  smooth repetition shaping, KL regularization, health-gated checkpoints, and
  exact deterministic resume.
  \item It redesigns the long-response execution path to avoid dense
  vocabulary-score retention, use rollout-only KV caching, recompute sampled
  token log-probabilities in bounded chunks, trim all-masked prompt columns,
  share the reference backbone, and split policy/value backward graphs.
  \item It evaluates Base, SFT, and PPO once on the same 2,017 prompts and
  publishes complete responses, reward scores, stopping and repetition
  metadata, resumable policy shards, review queues, and a public response
  explorer rather than reporting aggregate reward alone.
\end{enumerate}

The source repository is available at \projecturl{}. The response explorer is
available at \explorerurl{}.

\section{RLHF Formulation}

\subsection{Notation}

Let $x$ denote a formatted conversational prompt and
$y=(y_1,\ldots,y_T)$ an assistant response. The trainable policy is
$\pi_\theta(y\mid x)$, the frozen SFT reference is
$\pi_{\mathrm{ref}}(y\mid x)$, and the scalar reward model is
$r_\phi(x,y)$. A response mask $m_t\in\{0,1\}$ identifies generated assistant
tokens so prompt and padding positions do not contribute to policy or value
losses.

\subsection{Supervised fine-tuning}

SFT maximizes the likelihood of the preferred assistant response while
masking the prompt:

\begin{equation}
\mathcal{L}_{\mathrm{SFT}}(\theta)
=
-\sum_{t=1}^{T} m_t
\log \pi_\theta(y_t\mid x,y_{<t}).
\label{eq:sft}
\end{equation}

This is ordinary causal language modeling applied only to the completion. It
serves two roles: it creates a policy adapted to the preference dataset, and
it defines the reference distribution against which PPO drift is measured.
Low-rank adaptation (LoRA) \cite{hu2022} updates small trainable matrices
inside the attention and feed-forward projections while leaving the main
backbone fixed.

\subsection{Pairwise reward modeling}

Each preference record contains a chosen response $y^+$ and a rejected
response $y^-$ for the same prompt. The reward model is trained with the
Bradley--Terry logistic loss:

\begin{equation}
\mathcal{L}_{\mathrm{RM}}(\phi)
=
-\log \sigma\left(
r_\phi(x,y^+) - r_\phi(x,y^-)
\right),
\label{eq:rm}
\end{equation}

where $\sigma$ is the sigmoid function. Pairwise accuracy is

\begin{equation}
\mathrm{Acc}_{\mathrm{RM}}
=
\frac{1}{N}\sum_{i=1}^{N}
\mathbf{1}
\left[
r_\phi(x_i,y_i^+) > r_\phi(x_i,y_i^-)
\right].
\end{equation}

The sign of an isolated score is not intrinsically meaningful: adding a
constant to every reward leaves Equation~\ref{eq:rm} unchanged. This project
therefore interprets margins and rankings rather than treating zero as a
quality boundary. After training, a fixed demonstration-score offset is
estimated and subtracted consistently during PPO and evaluation.

\subsection{KL-regularized PPO}

For a sampled response, the reward model supplies a terminal score while the
reference policy supplies a token-level KL-shaped penalty:

\begin{align}
r_t^{\mathrm{KL}}
&=
-\beta\left[
\log \pi_\theta(y_t\mid x,y_{<t})
-
\log \pi_{\mathrm{ref}}(y_t\mid x,y_{<t})
\right],\\
R(x,y)
&=
r_\phi(x,y) + \sum_{t=1}^{T}r_t^{\mathrm{KL}}.
\label{eq:rlhf-reward}
\end{align}

The coefficient $\beta$ controls policy drift. Generalized advantage
estimation (GAE) computes

\begin{align}
\delta_t &= r_t + \gamma V(s_{t+1}) - V(s_t),\\
\hat{A}_t &= \delta_t + \gamma\lambda \hat{A}_{t+1},
\end{align}

with $\gamma=1$ and $\lambda=0.95$ in the final run. Advantages are whitened
over valid response tokens. For old-policy log probability
$\log\pi_{\mathrm{old}}$ and updated-policy log probability
$\log\pi_\theta$, define

\begin{equation}
\rho_t(\theta)
=
\exp\left(
\log\pi_\theta(y_t\mid s_t)
-
\log\pi_{\mathrm{old}}(y_t\mid s_t)
\right).
\end{equation}

PPO minimizes the negative clipped surrogate:

\begin{equation}
\mathcal{L}_{\mathrm{policy}}
=
-\mathbb{E}_t\left[
\min\left(
\rho_t\hat{A}_t,\,
\mathrm{clip}(\rho_t,1-\epsilon,1+\epsilon)\hat{A}_t
\right)
\right],
\label{eq:ppo}
\end{equation}

with $\epsilon=0.2$. The critic uses a clipped squared-error objective, and the
combined optimization loss is

\begin{equation}
\mathcal{L}
=
\mathcal{L}_{\mathrm{policy}}
+
c_v\mathcal{L}_{\mathrm{value}}
-
c_H\mathcal{H}(\pi_\theta).
\end{equation}

The final configuration uses $c_v=0.1$. These equations were implemented
directly in the custom trainer and later delegated to the experimental
\trl{} PPO trainer. The repository continues to own model preparation,
sampling controls, EOS handling, reward centering, checkpoint metadata, and
evaluation.

\section{Model, Data, and Preprocessing}

\subsection{Base model and preference data}

\dataset{} is a human-annotated preference dataset spanning general
instruction following, code, STEM, and multilingual tasks
\cite{wang2025}. Its breadth is useful for a general alignment study, but it
also makes the task harder than short-form summarization: many preferred
answers contain code, multi-step explanations, or long conversational context.

\begin{table}[H]
\centering
\caption{Final data and model profile.}
\label{tab:data-profile}
\begin{tabularx}{0.88\textwidth}{@{}Y r@{}}
\toprule
\textbf{Item} & \textbf{Value}\\
\midrule
Base policy & \model{}\\
Approximate parameter count & 0.5B\\
Raw training preference records & 38,459\\
Raw validation preference records & 2,017\\
Filtered SFT/RM training records & 36,264\\
Filtered SFT/RM validation pairs & 1,917\\
Unique PPO training prompts & 24,018\\
Policy-suite evaluation prompts & 2,017\\
Maximum SFT/RM sequence length & 4,096 tokens\\
Maximum PPO/evaluation prompt length & 3,072 tokens\\
PPO rollout response length & 768 tokens\\
Evaluation generation cap & 1,024 tokens\\
\bottomrule
\end{tabularx}
\end{table}

Invalid or tied preferences remove 2,195 training records and 100 validation
records from SFT/RM training. PPO uses prompt-only records and removes 14,441
duplicate training prompts. The policy-suite evaluation is separate from this
deduplicated PPO prompt cache and covers all 2,017 validation prompts.

\subsection{Chat formatting and truncation}

All stages use the Qwen chat template and the same tokenizer. The preprocessor
constructs three views of each retained preference record:

\begin{itemize}
  \item \textbf{SFT:} prompt plus chosen response, with a completion mask that
  excludes prompt tokens from Equation~\ref{eq:sft};
  \item \textbf{Reward model:} chosen and rejected sequences sharing the same
  prompt prefix;
  \item \textbf{PPO:} deduplicated, left-padded prompt token IDs.
\end{itemize}

PAD and EOS use distinct token IDs. Every SFT and reward-model completion ends
in EOS. When a sequence exceeds the budget, the preprocessor removes the
oldest non-system conversational turns first and left-truncates prompt tokens
only as a final fallback. It does not cut a response merely to preserve more
prompt history.

The project initially used a 1,024-token total limit. A direct length
diagnostic showed that this discarded too much preference signal:

\begin{table}[H]
\centering
\caption{Historical truncation diagnostic that motivated the 4,096-token
training budget.}
\label{tab:length-diagnostic}
\begin{tabular}{@{}rrrrr@{}}
\toprule
\textbf{Limit} &
\textbf{Train SFT} &
\textbf{Train RM} &
\textbf{Val. SFT} &
\textbf{Val. RM}\\
\midrule
1,024 & 38.47\% & 40.82\% & 36.78\% & 39.49\%\\
2,048 & 15.48\% & 16.47\% & 13.51\% & 14.87\%\\
3,072 & 5.28\% & 5.83\% & 4.69\% & 5.32\%\\
4,096 & 0.83\% & 1.00\% & 0.68\% & 0.89\%\\
\bottomrule
\end{tabular}
\end{table}

The final preparation report is consistent with this diagnostic: 304 of
36,264 SFT training prompts and 363 of 36,264 reward-model training prompts
required truncation. The median chosen/SFT sequence length is 768 tokens and
the 95th percentile is roughly 3,080 tokens. Long-context handling was
therefore not cosmetic; it preserved a substantial part of the dataset.

\section{Implementation and Experimental Infrastructure}

\subsection{From a custom trainer to TRL}

The first implementation built the algorithm from low-level components:
response masking, shifted token log-probabilities, token-level KL rewards,
GAE, normalized advantages, clipped policy and value losses, adaptive KL
control, LoRA checkpointing, and evaluation. This path was useful precisely
because it failed in visible ways. Short response caps produced truncation;
weak KL control or noisy adaptive estimates allowed drift; incorrect
checkpoint paths invalidated comparisons; and high reward sometimes
coincided with gibberish, empty outputs, or repeated text.

The final system uses Hugging Face \trl{} for trainer loops and standard
serialization \cite{trl}. The migration was not a deletion of the project.
Responsibilities were divided as follows:

\begin{table}[H]
\centering
\caption{Ownership after migration to \trl{}.}
\label{tab:ownership}
\begin{tabularx}{\textwidth}{@{}Y Y@{}}
\toprule
\textbf{Repository-owned} & \textbf{TRL-owned}\\
\midrule
HelpSteer3 parsing and filtering &
SFT, reward, and PPO trainer loops\\
Qwen chat formatting and truncation &
Mixed-precision/distributed preparation\\
Completion masks and prompt deduplication &
Gradient accumulation and optimizer stepping\\
Model initialization and reward centering &
PPO ratio, clipped objectives, and GAE\\
Fixed-length generation/EOS patches &
Trainer checkpoint serialization\\
Run manifests and resolved configurations &
Core training-loop metrics\\
Policy-suite evaluation and qualitative audit &
PEFT trainer integration\\
\bottomrule
\end{tabularx}
\end{table}

The custom implementation remains frozen at Git commit:
\begin{center}
{\small\texttt{6cbf214fcf1b91c7b756e303e533c2c86d2eba89}}
\end{center}
This preserves the learning history and provides a baseline for inspecting how
trainer mechanics changed.

The complete repository immediately before the later legacy-script pruning is
frozen at Git commit:
\begin{center}
{\small\texttt{6233219d2ee34517bc4203e0bb986f7cb27bf5d1}}
\end{center}
That snapshot retains every custom training entry point and historical helper.
The streamlined repository keeps only the active TRL SFT, reward, PPO, and DPO
path plus the shared evaluation, audit, tracking, and synchronization closure.
This second revision boundary makes removed educational code recoverable
without carrying it in the current working surface.

\subsection{Direct preference optimization extension}

The next controlled experiment uses Direct Preference Optimization (DPO)
\cite{rafailov2023} from the same frozen SFT checkpoint. This comparison
changes the alignment objective while holding the policy initialization,
preference source, tokenizer, sequence budget, LoRA parameterization, and
evaluation prompts fixed. It therefore tests whether direct offline preference
learning is more robust than optimizing the learned scalar reward with online
PPO at this model scale.

The DPO cache stores explicit \texttt{prompt}, \texttt{chosen}, and
\texttt{rejected} fields. It applies the same response-safe 4,096-token
preprocessing as SFT and reward modeling: complete EOS-terminated responses
are preserved, old prompt turns are removed first, and only a final prompt-side
token fallback is allowed. Tied pairs are excluded. The primary run uses one
row per non-tied pair because HelpSteer3's overall preference levels
$\{-3,\ldots,3\}$ are ordinal labels rather than established linear utility
distances. Preference magnitude, individual annotator scores, and directional
agreement remain attached as metadata for stratified analysis. A separately
named linear-replication run is available only as a sensitivity experiment.

The active DPO implementation uses the Hugging Face \trl{}
\texttt{DPOTrainer}, a standard sigmoid loss with $\beta=0.1$, zero dropout,
LoRA rank 16, a 4,096-token prepared-data budget, one epoch, and effective
batch size 32. Reference log probabilities are precomputed, and the reference
policy is the same SFT model with the trainable adapter disabled. On Colab,
training runs under \path{/content} while the SFT model, prepared DPO data,
periodic complete Trainer checkpoints, and final artifacts are synchronized
to Google Drive. Checkpoints are written every 200 optimizer updates and
\texttt{resume\_from\_checkpoint=auto} restores optimizer, scheduler, RNG, and
Trainer state after a runtime loss. The controlling notebook is
\path{notebooks/rlhf_dpo_colab_pipeline.ipynb}; the canonical configuration is
\path{configs/trl/qwen25_05b_helpsteer3_dpo.yaml}.

\subsection{Implementation details adopted from N+}

RLHF is unusually sensitive to details that can look incidental in pseudocode.
The N+ reproduction \cite{huang2024}, which itself reconstructs the
summarization setup of Stiennon et al.\ \cite{stiennon2020}, guided the final
configuration:

\begin{table}[H]
\centering
\small
\caption{High-impact RLHF/PPO implementation details in the final path.}
\label{tab:nplus}
\begin{tabularx}{\textwidth}{@{}p{0.28\textwidth}Y@{}}
\toprule
\textbf{Detail} & \textbf{Implementation}\\
\midrule
Disable PPO dropout &
PyTorch dropout probabilities and LoRA dropout are zero, preventing stochastic
ratio noise between rollout and update.\\
Match behavior probabilities &
Generation temperature is applied consistently when storing and recomputing
log probabilities; checkpoint generation heuristics are neutralized.\\
Preserve completions and EOS &
Prompt history is truncated before response tokens; SFT/RM completions end in
EOS.\\
Initialize RM from SFT &
The merged SFT model supplies the reward backbone.\\
Initialize reward head deliberately &
The scalar head uses zero bias and standard deviation
$1/\sqrt{d_{\mathrm{hidden}}+1}$.\\
Center rewards &
Training uses a centering regularizer, then subtracts a persisted mean
demonstration score.\\
Initialize critic from RM &
The PPO value model begins from the trained reward-model weights.\\
Use a frozen SFT reference &
PPO adapter weights are disabled for reference log probabilities.\\
Fixed-length EOS trick &
Rollouts sample the configured length without stopping at EOS, are truncated
at EOS afterward. The evaluated historical run used reward $-1$ when EOS was
missing; the guarded follow-up instead replaces the score with the calibrated
0.5th-percentile lower reward boundary.\\
Whiten rewards/advantages &
The guarded follow-up disables optional reward whitening so calibrated terminal
penalties retain their scale; \trl{} continues to whiten masked advantages.\\
Adam numerical setting &
Adam epsilon is $10^{-5}$.\\
\bottomrule
\end{tabularx}
\end{table}

These choices reduce avoidable instability; they cannot repair a weak reward
model, noisy preferences, or insufficient model capacity.

\subsection{Reproducibility and Colab execution}

Each stage writes a resolved YAML configuration, experiment manifest, summary,
trainer metrics, and stage-specific artifacts. Smoke runs precede expensive
runs. A ``doctor'' script validates the Python interpreter, CUDA visibility,
trainer APIs, prepared datasets, tokenizer special tokens, and local/Drive
write paths.

The Colab workflow keeps active datasets and checkpoints on local SSD for
throughput and synchronizes periodic checkpoints and final artifacts to Google
Drive. SFT and reward training support exact Transformers checkpoint resume.
The guarded PPO wrapper adds an explicit exact-resume layer around the
experimental \trl{} trainer. Every verified checkpoint contains the LoRA policy,
complete value model, optimizer, scheduler, Trainer state, Python/NumPy/CPU/CUDA
RNG state, deterministic rollout-plan position, calibrated guardrails, and a
health marker written only after all required files exist. Training is divided
into cumulative target updates and resumes from the latest verified checkpoint;
an incomplete newer directory is ignored.

Evaluation is resumable at the policy-output level. Each policy writes
complete JSONL records as generation proceeds; ``partial'' denotes
interrupt-safe stage output, not truncated response text. A fingerprint guards
against reusing cached responses under a changed checkpoint or generation
configuration.

\section{Final Training Configuration and Results}

\subsection{Supervised fine-tuning}

SFT trains LoRA modules in the attention projections
(\texttt{q}, \texttt{k}, \texttt{v}, and \texttt{o}) and feed-forward
projections (\texttt{gate}, \texttt{up}, and \texttt{down}). Prompt and
padding tokens are masked from the loss, so the objective applies only to the
preferred assistant response. The effective batch size is $8\times4=32$ on
one device.

\begin{table}[H]
\centering
\caption{Final SFT configuration and outcome.}
\label{tab:sft}
\begin{tabularx}{0.88\textwidth}{@{}Y r@{}}
\toprule
\textbf{Setting or metric} & \textbf{Value}\\
\midrule
Trainer & \trl{} \texttt{SFTTrainer}\\
Epochs & 1\\
LoRA rank / alpha & 16 / 32\\
LoRA dropout & 0\\
Effective batch size & 32\\
Learning rate & $5\times10^{-6}$\\
Scheduler & cosine, 3\% warmup\\
Precision & bfloat16, TF32 enabled\\
Maximum total length & 4,096 tokens\\
Train / validation loss & 1.0556 / 1.1127\\
Validation mean token accuracy & 72.02\%\\
\bottomrule
\end{tabularx}
\end{table}

The merged SFT checkpoint is used in three roles: as a comparison policy, as
the initialization for the reward model and PPO policy, and as the frozen PPO
reference. That common starting point makes downstream changes easier to
attribute.

\subsection{Reward-model training}

The reward model starts from the merged SFT backbone and trains LoRA rank-32
adapters plus a scalar head using the Bradley--Terry loss. A persisted scalar
offset centers preferred demonstrations consistently in PPO and evaluation;
it does not change pairwise ordering.

\begin{table}[H]
\centering
\caption{Final two-epoch reward-model audit.}
\label{tab:rm}
\begin{tabularx}{0.88\textwidth}{@{}Y r@{}}
\toprule
\textbf{Metric} & \textbf{Value}\\
\midrule
Usable validation pairs & 1,917\\
Validation loss & 0.6166\\
Pairwise accuracy & 65.62\%\\
Mean chosen--rejected margin & 0.3578\\
Persisted score offset & $-0.1985$\\
Code / General accuracy & 70.23\% / 62.91\%\\
STEM / Multilingual accuracy & 59.26\% / 71.82\%\\
\bottomrule
\end{tabularx}
\end{table}

The domain spread matters for interpretation. The model is useful as a broad
preference proxy, but 65.62\% accuracy---and especially the weaker STEM and
general subsets---is not sufficient to treat its scalar score as ground truth.
Only score differences are identifiable; negative absolute scores do not by
themselves indicate a bad response.

\subsection{Final guarded PPO run}

The selected PPO run was designed around long-form responses rather than around
maximizing the reward model at any cost. It completed the configured 12,000-
episode budget as 12,032 rollouts: 188 batches of 64 responses. Each batch
contains 16 code, 16 general, 16 STEM, and 16 multilingual prompts. Approximately
12,007 prompt rows are unique, so the run covers a broad slice of the dataset
without pretending to be a conventional supervised epoch.

The terminal score is clipped to empirical 0.5th and 99.5th percentiles
estimated from 4,096 stratified reward-training pairs,
$[-2.7507,1.5018]$. A response that fails to produce EOS receives the lower
bound, which makes nontermination as unattractive as the worst calibrated
response rather than assigning an arbitrary weak penalty. A smooth repeated
word-level 4-gram penalty activates only above 0.3084, the preferred-response
95th percentile. This protects ordinary repeated terminology while shaping
clear loops. The PPO objective retains a frozen SFT reference, KL coefficient
0.10, clipped policy and value losses, an RM-initialized critic, advantage
normalization, and four PPO epochs per rollout batch. Scalar reward whitening
is disabled so the calibrated reward scale remains meaningful.

\begin{table}[H]
\centering
\caption{Final guarded PPO configuration and terminal diagnostics.}
\label{tab:ppo-final}
\begin{tabularx}{0.94\textwidth}{@{}Y r@{}}
\toprule
\textbf{Setting or metric} & \textbf{Value}\\
\midrule
Prompt / response limits & 3,072 / 768 tokens\\
Completed rollouts / updates & 12,032 / 188\\
Rollout batch / PPO epochs & 64 / 4\\
Domain composition per update & 16 / 16 / 16 / 16\\
Learning rate & $3\times10^{-6}$, linear decay\\
KL coefficient & 0.10\\
Calibrated reward interval & $[-2.7507,1.5018]$\\
Repetition activation threshold & 0.3084\\
Final objective KL to SFT reference & 0.7305\\
Final objective KL per response token & 0.00199\\
Final old-to-new approximate KL & 0.000537\\
Final clip fraction & 0.00571\\
Final value loss & 0.1076\\
Final rollout mean / median tokens & 366.84 / 342.5\\
Final rollout EOS / cap rate & 82.81\% / 17.19\%\\
Final rollout repeated-4-gram fraction & 17.0\%\\
Empty-response rate & 0\%\\
Peak allocated CUDA memory & 37.64 GiB\\
\bottomrule
\end{tabularx}
\end{table}

Across updates 1--94 and 95--188, mean reward-model score improves from
$-1.110$ to $-0.961$, EOS rate rises from 74.3\% to 81.0\%, cap rate falls
from 25.7\% to 19.0\%, mean response length falls from 452 to 409 tokens, and
mean repeated-4-gram fraction falls from 22.2\% to 19.5\%. Mean value loss
falls from 2.83 to 0.114 as the critic adapts to the reward scale. These trends
support a narrow statement: the run learned and remained numerically healthy.
They do not establish that the learned reward is an independent measure of
human preference.

\subsection{Memory-efficient long-response execution}

Long responses made systems design part of the PPO method. Stock experimental
\trl{} retained a dense generation-score tensor with shape
$64\times768\times151{,}665$, approximately 27.77 GiB in float32 before the
policy update. The final execution path instead:

\begin{itemize}
  \item generates token IDs without retaining per-step vocabulary scores;
  \item uses a rollout-only KV cache, then recomputes behavior log-probabilities
  only for sampled response tokens;
  \item bounds the recomputation with \texttt{logits\_to\_keep=769} and
  response-position chunking;
  \item dynamically drops prompt columns masked for every item in the batch;
  \item shares the reference backbone by evaluating with the policy adapter
  disabled rather than loading a second full backbone; and
  \item backpropagates policy and value losses through separate graphs before
  one optimizer step, reducing peak activation retention.
\end{itemize}

These changes preserve the rollout distribution, reward, generalized
advantages, clipping, domain balance, and effective batch size. They change
only how equivalent quantities are materialized. Together with bfloat16 and
TF32, they make 768-token, batch-64 PPO practical with 37.64 GiB peak allocated
memory and no generation or log-probability OOM fallback in the completed run.

\section{Experimental Progression}

The project did not move directly from a first implementation to the final
run. Table~\ref{tab:history} records the main stages; detailed custom metrics
are retained in Appendix A.

\begin{table}[H]
\centering
\small
\caption{Selected experimental progression. Only the last row is the final
reported PPO result.}
\label{tab:history}
\begin{tabularx}{\textwidth}{@{}p{0.18\textwidth}p{0.17\textwidth}Y@{}}
\toprule
\textbf{Stage} & \textbf{Budget} & \textbf{What it established}\\
\midrule
Short custom smoke tests & 96--128 output tokens &
Validated masks, KL, value learning, and checkpointing, while exposing
truncated examples, gibberish, empty responses, and invalid checkpoint loading.\\
Custom long-512 PPO & 397 updates, batch 2 &
Produced complete responses but weak aggregate preference results; motivated
full-suite evaluation and stronger stopping diagnostics.\\
Custom 1,024-token evaluation & 2,017 prompts &
Reduced truncation but exposed high-reward loops, fabricated details, and
reward-model mismatches; not a clean cap ablation.\\
TRL checkpoint 100 & 6,400 rollouts, 100 updates &
Reached 50.92\% proxy win rate versus Base and 57.71\% versus SFT, but with
520-token median responses, 27.42\% cap hits, and 31.88\% above the 25\%
repetition threshold. It demonstrated stronger proxy optimization and clearer
reward exploitation, not a final aligned policy.\\
Final guarded TRL PPO & 12,032 rollouts, 188 updates &
Completed the declared budget with calibrated rewards, EOS and repetition
guards, balanced domains, exact resume, and the memory-efficient long-response
path. This is the final reported policy.\\
\bottomrule
\end{tabularx}
\end{table}

The checkpoint-100 experiment remains important. It is the clearest evidence
that the learned reward can be improved while visible behavior deteriorates.
Relative to that archived run, the final guarded evaluation is shorter and
less capped and repetitive, but it also loses the earlier reward-model lead.
Because generation caps and other execution details differ, the comparison is
an engineering progression rather than a controlled causal ablation.

\section{Final Policy-Suite Evaluation}

\subsection{Protocol}

Base, SFT, and guarded PPO each generate one response for every one of the
2,017 validation prompts. Generation uses the same tokenizer, prompt formatting,
3,072-token prompt limit, 768-token response limit, temperature 0.7, and top-p
1.0. The two-epoch reward model scores every response. Evaluation writes one
JSONL shard per policy as it proceeds and can resume without regenerating
completed rows. All pairwise comparisons are derived from the same shared
response table, eliminating inconsistent prompt samples across comparisons.

The reward-model winner is useful but incomplete. The report therefore places
length, cap, EOS, and repeated-4-gram diagnostics before the win statistics and
retains full responses for human inspection.

\subsection{Overall results}

\begin{table}[H]
\centering
\small
\caption{Behavior-first full-suite summary on 2,017 prompts.}
\label{tab:overall}
\begin{tabular}{@{}lrrrrrrrr@{}}
\toprule
 & \multicolumn{5}{c}{\textbf{Behavior}} & \multicolumn{3}{c}{\textbf{Learned reward}}\\
\cmidrule(lr){2-6}\cmidrule(l){7-9}
\textbf{Policy} & \textbf{Mean tok.} & \textbf{Med. tok.} & \textbf{Cap} &
\textbf{Rep. $>25$} & \textbf{Rep. $>50$} & \textbf{Mean} &
\textbf{Wins} & \textbf{Rate}\\
\midrule
Base & 373.5 & 333 & 16.21\% & 165 (8.18\%) & 43 (2.13\%) & 0.2229 & 845 & 41.89\%\\
SFT & 392.5 & 364 & 18.34\% & 374 (18.54\%) & 143 (7.09\%) & 0.2037 & 585 & 29.00\%\\
PPO & 384.1 & 355 & 16.31\% & 326 (16.16\%) & 105 (5.21\%) & 0.1753 & 563 & 27.91\%\\
Tie & -- & -- & -- & -- & -- & -- & 24 & 1.19\%\\
\bottomrule
\end{tabular}
\end{table}

\begin{table}[H]
\centering
\caption{Pairwise reward-model comparisons.}
\label{tab:pairwise}
\begin{tabular}{@{}lrrrrr@{}}
\toprule
\textbf{Comparison} & \textbf{Left} & \textbf{Right} & \textbf{Ties} &
\textbf{Right rate} & \textbf{Mean $\Delta$}\\
\midrule
Base vs SFT & 1,145 & 849 & 23 & 42.09\% & $-0.0192$\\
Base vs PPO & 1,179 & 817 & 21 & 40.51\% & $-0.0476$\\
SFT vs PPO & 1,018 & 974 & 25 & 48.29\% & $-0.0284$\\
\bottomrule
\end{tabular}
\end{table}

PPO is closest to SFT, but the reward proxy does not prefer it in aggregate.
The mean deltas are small and the win counts still favor Base and SFT. This is
not a hidden victory and should not be advertised as one. It is evidence that
the policy remained competitive while the guardrails constrained several
behaviors that the earlier checkpoint had exploited.

\subsection{Domain behavior}

\begin{table}[H]
\centering
\caption{Guarded PPO versus Base by prompt domain.}
\label{tab:domains}
\begin{tabular}{@{}lrrrrr@{}}
\toprule
\textbf{Domain} & \textbf{Prompts} & \textbf{PPO wins} &
\textbf{Base wins} & \textbf{Ties} & \textbf{PPO rate}\\
\midrule
Code & 438 & 157 & 279 & 2 & 35.84\%\\
General & 931 & 400 & 517 & 14 & 42.96\%\\
STEM & 245 & 96 & 148 & 1 & 39.18\%\\
Multilingual & 403 & 164 & 235 & 4 & 40.69\%\\
\bottomrule
\end{tabular}
\end{table}

PPO does not beat Base within any domain. Against SFT, however, it wins the
STEM subset 133 to 109 with three ties. This reinforces the need to inspect
both the comparison policy and the task domain rather than collapse the whole
project into one scalar win rate.

\section{Qualitative Audit and Response Explorer}
\label{sec:audit}

For response word-level 4-gram multiset $\mathcal{G}$, the repeated fraction is
\begin{equation}
f_{\mathrm{rep4}}=1-\frac{|\mathrm{unique}(\mathcal{G})|}{|\mathcal{G}|}.
\end{equation}
It is a review heuristic, not a semantic metric: legitimate code and technical
prose repeat terms, while some loops change punctuation enough to evade it.

The final audit creates overlapping navigation queues rather than mutually
exclusive quality labels: 9 qualified PPO candidates, 120 strong PPO-loss
candidates, 326 repetition-risk rows, and 155 high-reward repetition or
stopping mismatches. The public explorer exposes all 2,017 prompts, complete
Base and PPO responses, reward scores and margins, token counts, EOS/cap state,
repetition diagnostics, domain filters, and a balanced curated first pass.
Reward scores are explicitly labeled as diagnostics, not human judgments.

Manual inspection finds useful local changes: clearer structure, better
coverage of some multi-part instructions, supportive framing, and more compact
answers in selected cases. It also finds failures that the learned reward does
not reliably detect: invalid code, confused factual or scientific reasoning,
prompt restatement, fabricated details or citations, irrelevant continuation,
and occasional near-cap loops that still receive high reward. These examples
are the practical evidence for reward-model overoptimization discussed more
generally by Gao et al.\ \cite{gao2023}.

The explorer is therefore not a promotional gallery of PPO wins. It is the
bridge from aggregate statistics to the actual object being optimized: model
text. The repository retains both positive and negative queues so a reader can
inspect where the proxy agrees with visible quality and where it does not.

\section{Discussion}

The completed run supports a strong systems conclusion and a modest alignment
conclusion. The system joins long-context preprocessing, response-only SFT,
pairwise reward learning, online PPO, calibrated safeguards, exact resume,
memory-bounded long-response execution, full-suite generation, and qualitative
audit on a heterogeneous preference dataset. The optimizer behaves as intended:
losses remain finite, old-to-new approximate KL and clip fraction are small,
EOS behavior improves during training, empty outputs remain absent, and the
critic converges to the reward scale.

That stability is not equivalent to model improvement. The same imperfect
reward model supplies the training reward and the final quantitative judge;
there is no blinded multi-rater human study, no independent frontier-model
judge, and only one sampled generation per policy and prompt. Response length
remains a confound, the 0.5B policy and reward model have limited capacity, and
word-level repetition heuristics cannot establish factuality, safety, or
instruction following. Historical token-budget and checkpoint comparisons also
change more than one variable and should be read as operational evidence, not
clean ablations.

The result is consequently framed without hiding the negative evidence. PPO
changed behavior and produced useful examples, but guarded PPO did not beat
Base or SFT under the learned proxy. The project succeeds because it makes
that outcome diagnosable and reproducible, and because the earlier higher-proxy
checkpoint is preserved as an explicit example of why reward optimization alone
is not enough.

\section{Future Work}

The highest-value next step is independent evaluation rather than another
unqualified PPO continuation. A blinded human or strong-LLM comparison on a
stratified sample should measure agreement with the learned reward model and
identify hard negatives for reward-model retraining. Code responses should be
executed where possible; STEM and factual claims should receive task-specific
checks. A controlled decoding study can then compare 512, 768, and 1,024-token
budgets while holding the checkpoint, prompt order, seed, batch size, and
software environment fixed.

Only after the judge is stronger should PPO be extended or restarted. Useful
variants include a response-length curriculum, checkpoint selection based on a
joint scorecard of reward, KL, EOS, cap, repetition, and independent preference,
and larger reward/policy models. The prepared DPO implementation provides an
offline preference baseline from the same SFT checkpoint and data, but it is an
archived extension until it is actually trained and evaluated under the same
suite.

\section{Conclusion}

This project completes and audits an end-to-end RLHF pipeline for a small
instruction model on heterogeneous, long-form preference data. Long-context SFT
retains nearly all responses; the reward model learns a usable but imperfect
pairwise signal; guarded PPO completes 12,032 rollouts with stable optimization,
health-gated resume, explicit anti-exploitation controls, and a practical
long-response memory path. The full evaluation does not show a PPO win over
Base. It shows something more useful for this project: a controlled policy that
can be inspected alongside the proxy failures and qualitative tradeoffs that
shaped it.

The durable result is therefore not a single reward number. It is the combined
training, systems, evaluation, and audit record---including the runs that looked
better to the reward model but worse on stopping and repetition. That record is
a defensible conclusion for this phase and a concrete foundation for future
alignment experiments.

\appendix

\section{Detailed Historical Experiment Ledger}

This appendix retains the operational facts behind the compact progression in
Table~\ref{tab:history}. Historical paths and configurations are evidence of
what was run; they are not active recommendations.

\subsection{Initial validation and sequence-length diagnosis}

Early 96- and 128-token generations accelerated debugging but often stopped
inside code blocks or explanations. These runs exposed multilingual or
gibberish drift, vulgar debug-pattern text, EOS and blank-response collapse,
overaggressive reward shaping, and invalid comparisons caused by incorrect
checkpoint paths. The resulting safeguards included prompt sanitation, KL
anchoring, empty-response checks, repeated-text diagnostics, and explicit
checkpoint validation.

The sequence-length diagnostic established why the project moved from 1,024
to 4,096 total tokens:

\begin{table}[H]
\centering
\caption{Fraction of records requiring truncation at each total-token limit.}
\label{tab:length-ledger}
\begin{tabular}{@{}rrrrr@{}}
\toprule
\textbf{Limit} & \textbf{Train SFT} & \textbf{Train RM} &
\textbf{Valid. SFT} & \textbf{Valid. RM}\\
\midrule
1,024 & 38.47\% & 40.82\% & 36.78\% & 39.49\%\\
2,048 & 15.48\% & 16.47\% & 13.51\% & 14.87\%\\
3,072 & 5.28\% & 5.83\% & 4.69\% & 5.32\%\\
4,096 & 0.83\% & 1.00\% & 0.68\% & 0.89\%\\
\bottomrule
\end{tabular}
\end{table}

\subsection{Frozen custom SFT, reward, and PPO baseline}

The pre-TRL lineage is frozen at commit
\texttt{6cbf214fcf1b91c7b756e303e533c2c86d2eba89}. Custom SFT used
\model{}, 4,096 total tokens, two epochs, batch size 6, accumulation 3,
learning rate $5\times10^{-6}$, and LoRA rank 16. It completed 4,030 optimizer
steps and wrote to
\path{outputs/rlhf/qwen25_05b_helpsteer3_sft_4096/}.

The two-epoch custom reward model wrote to
\path{outputs/rlhf/qwen25_05b_helpsteer3_reward_4096_epoch2/} and produced the
following held-out audit:

\begin{table}[H]
\centering
\caption{Historical custom reward-model result.}
\begin{tabular}{@{}lr@{}}
\toprule
\textbf{Metric} & \textbf{Value}\\
\midrule
Validation pairs & 1,917\\
Accuracy / loss & 71.62\% / 0.9734\\
Mean chosen--rejected margin & 0.9094\\
Code / General accuracy & 74.88\% / 71.01\%\\
STEM / Multilingual accuracy & 63.37\% / 75.15\%\\
\bottomrule
\end{tabular}
\end{table}

The custom long-512 PPO run used a 3,072-token prompt budget, 512-token
rollouts, learning rate $3\times10^{-7}$, policy clip 0.06, one PPO epoch,
rollout batch 2, and an adaptive KL coefficient initialized at 0.18 with a
0.14 minimum. It completed 397 of 400 requested updates and ended at KL
coefficient 0.14. Rewards were clipped to $[-5,5]$; scalar reward whitening
was absent; length penalty was 0.002; missing-EOS penalty was zero; responses
shorter than 32 tokens received a 0.5 penalty; group-relative normalization
was disabled. Generation applied repetition penalty 1.05 and blocked repeated
4-grams. Prompt batches were left padded, and response masks excluded padding
from PPO losses. Its output is
\path{outputs/rlhf/qwen25_05b_helpsteer3_ppo_4096_epoch2_long512/}.

\subsection{Historical policy-suite results}

The 512-token suite used all 2,017 validation prompts. Three-way wins were
Base 827 (41.00\%), SFT 556 (27.57\%), PPO 525 (26.03\%), and 109 ties
(5.40\%). Mean rewards were $-3.5339$, $-3.4280$, and $-3.6666$; median
response lengths were 332, 363, and 356 tokens; and cap-hit rates were
29.90\%, 29.95\%, and 29.70\%. Its complete pairwise table is retained below.

\begin{table}[H]
\centering
\caption{Historical custom 512-token pairwise evaluation.}
\begin{tabular}{@{}lrrrrr@{}}
\toprule
\textbf{Comparison} & \textbf{Left} & \textbf{Right} & \textbf{Ties} &
\textbf{Right rate} & \textbf{Mean $\Delta$}\\
\midrule
Base vs SFT & 1,044 & 927 & 46 & 45.96\% & $+0.1059$\\
Base vs PPO & 1,068 & 904 & 45 & 44.82\% & $-0.1327$\\
SFT vs PPO & 965 & 898 & 154 & 44.52\% & $-0.2386$\\
\bottomrule
\end{tabular}
\end{table}

The later 1,024-token suite produced three-way wins of Base 978 (48.49\%),
SFT 475 (23.55\%), PPO 467 (23.15\%), and 97 ties (4.81\%). Mean rewards were
$-3.3634$, $-3.6114$, and $-3.5771$; median lengths were 334, 360, and 363;
cap-hit rates were 8.08\%, 10.16\%, and 11.60\%; empty rates were zero.

\begin{table}[H]
\centering
\caption{Historical custom 1,024-token pairwise evaluation.}
\begin{tabular}{@{}lrrrrr@{}}
\toprule
\textbf{Comparison} & \textbf{Left} & \textbf{Right} & \textbf{Ties} &
\textbf{Right rate} & \textbf{Mean $\Delta$}\\
\midrule
Base vs SFT & 1,215 & 763 & 39 & 37.83\% & $-0.2480$\\
Base vs PPO & 1,190 & 785 & 42 & 38.92\% & $-0.2137$\\
SFT vs PPO & 963 & 892 & 162 & 44.22\% & $+0.0343$\\
\bottomrule
\end{tabular}
\end{table}

PPO's 785 Base wins had mean margin $+1.6210$, while its 1,190 losses had
mean margin $-1.4315$. Against SFT, its 892 wins averaged $+1.3503$ and its
963 losses $-1.1789$. Domain PPO win rates against Base were 26.26\% for code
(115/438), 44.47\% for general (414 wins, 487 losses, 30 ties), 39.59\% for
STEM (97/245 with two ties), and 39.45\% for multilingual (159 wins, 236
losses, eight ties).

Increasing the evaluation cap reduced Base/SFT/PPO cap hits from 603/604/599
to 163/205/234. It also exposed more PPO repetition: 224 responses exceeded
25\% repeated 4-grams in the 512-token suite versus 325 (16.11\%) in the
1,024-token suite; severe cases above 50\% rose from 74 to 156. Manual review
found fabricated citations, incorrect chemistry, irrelevant continuations,
and high-reward loops. This was not a clean token-cap ablation: batch size
changed from 8 to 128, and bfloat16 batching, environment, or evaluator
revisions could change close greedy continuations.

\subsection{Historical A100 40GB profile}

An earlier speed-oriented Colab profile used full bfloat16 weights rather than
4-bit loading. Reward training used batch 8, accumulation 2 (effective 16),
all filtered HelpSteer3 pairs, and refreshed token throughput, CUDA memory,
CSV files, and plots every configured artifact interval. Batch 12 was the
suggested increase below roughly 28--30GB peak memory; batch 6 was the fallback
near 36GB or after OOM.

Historical PPO used full bfloat16 policy/reference/reward models, rollout batch
16, minibatch 4, prompt length 768, response length 160, checkpoints every 50
updates, and metric refresh every 25 updates. Rollout batch 24 was the first
recommended utilization increase below 30--32GB; minibatch 6 or 8 was reserved
for runs with measured headroom. The early evaluator covered 200 prompts;
the completed suite superseded it with all 2,017.

\section{Experiment Tracking and Run Governance}

The project uses a filesystem contract rather than a tracking server. Every
new run directory contains the following where applicable:

\begin{table}[H]
\centering
\caption{Run-tracking contract.}
\begin{tabularx}{0.94\textwidth}{@{}p{0.31\textwidth}Y@{}}
\toprule
\textbf{Artifact} & \textbf{Purpose}\\
\midrule
\texttt{config\_resolved.yaml} & Exact configuration consumed by the run.\\
\texttt{experiment\_manifest.json} & Identity, intent, config hash, command,
source state, environment, attempts, status, and summary.\\
\texttt{run\_metadata.json} & Backward-compatible hardware/software record.\\
\texttt{*\_metrics.jsonl} / CSV & Step-level measurements.\\
Run or evaluation summary & Final aggregate result.\\
Checkpoints, samples, plots & Stage-specific evidence.\\
\bottomrule
\end{tabularx}
\end{table}

An interrupted process leaves its manifest in \texttt{running} state.
Reusing the directory appends an attempt and records whether the resolved
configuration changed; successful completion updates status and summary.
Configs can declare non-behavioral intent---experiment ID, name, group, tags,
hypothesis, and parent runs---while algorithmic choices remain in functional
sections such as generation, reward shaping, KL, and PPO.

Runs are compared with
\texttt{python -m scripts.rlhf\_compare\_runs RUN\_A RUN\_B}; metadata-only
fields are hidden by default. The command supports \texttt{--include-metadata},
\texttt{--format json}, and \texttt{--output comparison.md}. The
frozen custom baseline remains under
\path{experiments/baselines/qwen25_05b_helpsteer3_ppo_long512/}; its YAML
record is authoritative after removal of its redundant prose README.

The final records live under
\path{rlhf_runs/qwen25_05b_helpsteer3_trl_a100_full/full/}. Stage sources of
truth are the SFT \texttt{run\_summary.json}, reward epoch-1 and epoch-2
summaries and audits, the guarded PPO \texttt{run\_summary.json}, the
health-verified \texttt{checkpoint-188}, and the final 2,017-row policy-suite
summary, sample table, and audit queues. Historical checkpoint-100 artifacts
remain available only for the experimental progression and are not the final
policy. Persisted resolved configuration and the executed guarded PPO notebook
take precedence over obsolete labels in older directories.

\section{Operational TRL, Colab, and Evaluation Contract}

The stage datasets share one filtered source but have distinct contracts: SFT
stores token IDs and a completion mask; reward modeling stores paired chosen
and rejected token IDs with a common prefix; PPO stores deduplicated
left-padded prompts; DPO stores explicit prompt/chosen/rejected text plus
preference-strength and annotator metadata. Prompt history is shortened before
any final prompt-side truncation, responses retain EOS, and PAD is distinct
from EOS.

The preflight doctor validates the exact interpreter, package versions, CUDA,
trainer APIs, tokenizer special tokens, prepared data, and local/Drive write
permissions before expensive work. Evaluation always loads the SFT tokenizer;
Base is resized once for its distinct PAD token. Neutral generation controls
are explicit so checkpoint defaults cannot create asymmetric decoding. Policy
list overrides accept either dotted or bracketed indices. The policy labeled
Base must keep a null checkpoint and load the configured model name; pointing
it at SFT would silently make the Base and SFT columns identical.

Active work uses Colab local SSD for model shards, datasets, logs, and Trainer
state, with periodic checkpoints and final artifacts copied to mounted Drive.
This is faster and keeps model tensors outside Git. SFT, reward, and DPO
support exact Trainer resume. DPO's \texttt{auto} mode restores model,
optimizer, scheduler, RNG, and Trainer state. The guarded wrapper supplies the
missing protocol around experimental PPO: policy/value state, optimizer,
scheduler, Trainer state, RNG state, rollout position, and guardrail
fingerprints must all pass a health check before a directory is eligible for
automatic continuation.

Evaluation is resumable per policy. Its ``partial policy outputs'' are complete
JSONL rows written incrementally; \emph{partial} describes interrupt-safe run
completion, not shortened prompts or responses. A fingerprint prevents cached
rows from being reused when checkpoint or generation settings change.

TRL PPO APIs are experimental, and the repository intentionally validates the
installed stack instead of assuming that an unpinned upgrade is compatible.
Every dependency change requires a smoke run and a recorded environment before
an expensive experiment.

\section{Curation Workflow and Publication Standard}

The final workflow treats curation as review navigation rather than a source of
headline metrics. The following artifacts are generated directly from the
2,017-row shared policy table:

\begin{table}[H]
\centering
\small
\caption{Curation artifacts retained by the final workflow.}
\begin{tabularx}{\textwidth}{@{}p{0.43\textwidth}Y@{}}
\toprule
\textbf{Artifact} & \textbf{Purpose}\\
\midrule
\texttt{policy\_suite\_samples.jsonl} / CSV / XLSX & Full prompts,
Base/SFT/PPO responses, learned rewards, lengths, EOS/cap state, and pairwise
deltas.\\
\texttt{qualitative\_audit\_auto.md} & Aggregate stopping/repetition audit and
candidate summary.\\
\texttt{curation\_qualified\_ppo\_candidates.csv} & Low-warning PPO candidates
for positive review.\\
\texttt{curation\_strong\_ppo\_losses.csv} & Large learned-reward regressions.\\
\texttt{curation\_ppo\_repetition\_risks.csv} & Responses above the configured
repeated-4-gram review threshold.\\
\texttt{curation\_reward\_model\_mismatches.csv} & High learned reward combined
with repetition or stopping warnings.\\
\path{portfolio_full_policy_comparisons_final_trl.json} & Redacted full
response-explorer payload.\\
\path{portfolio_curated_100_manifest_final_trl.json} & Balanced public
first-pass subset.\\
\bottomrule
\end{tabularx}
\end{table}

The queues overlap and are deterministic functions of recorded diagnostics;
they are not manual, human, or LLM-as-a-judge labels. Review must cover apparent
wins, near ties, strong losses, repetition, proxy mismatches, correctness-sensitive
code and STEM cases, and multilingual drift. Before calling an example an
improvement, a reviewer checks instruction following, correctness, completeness,
relevance, repetition, safety, and whether the reward advantage reflects
substance rather than length or formatting. Publication must show both gains
and failures. The defensible conclusion is that guarded PPO is stable and
inspectable but does not beat Base under the learned judge.

\section{Reproducibility and Artifact Map}

\small
\begin{longtable}{@{}p{0.27\textwidth}p{0.67\textwidth}@{}}
\caption{Primary project artifacts used in this report.}
\label{tab:artifacts}\\
\toprule
\textbf{Artifact} & \textbf{Purpose}\\
\midrule
\endfirsthead
\toprule
\textbf{Artifact} & \textbf{Purpose}\\
\midrule
\endhead
Project overview &
\path{README.md}. Concise final guarded-run story, training profile, full-suite
results, interpretation, and links.\\
Repository guide &
\path{docs/REPOSITORY_GUIDE.md}. Installation, commands, directory map,
configuration conventions, checkpoint semantics, evaluation resume, and the
archived DPO extension.\\
Academic report &
\path{docs/rlhf_project_report.tex} and PDF. Concise scientific account of SFT,
reward modeling, guarded PPO, memory execution, and the final evaluation.\\
Technical record &
\path{docs/rlhf_technical_companion.tex} and PDF. Chronological experiments,
trainer migration, earlier PPO results, operational contracts, DPO design,
audit, and artifact governance.\\
Guarded PPO notebook &
\path{notebooks/rlhf_ppo_guarded_colab_pipeline.ipynb}. Executed 188-update PPO
workflow and synchronization steps.\\
Evaluation notebook &
\path{notebooks/rlhf_full_eval_and_curation.ipynb}. Full-suite analysis and
curation workflow.\\
Final SFT summary &
\path{rlhf_runs/qwen25_05b_helpsteer3_trl_a100_full/full/sft/run_summary.json}.\\
Final reward audit &
\path{rlhf_runs/qwen25_05b_helpsteer3_trl_a100_full/full/reward_epoch2/reward_audit.json}.\\
Final PPO summary &
\path{rlhf_runs/qwen25_05b_helpsteer3_trl_a100_full/full/ppo_guarded_r768_b64_kl10_rm_ep2/run_summary.json}.\\
Final checkpoint &
\path{.../ppo_guarded_r768_b64_kl10_rm_ep2/checkpoint-188/}. LoRA policy,
value model, optimizer, scheduler, Trainer state, RNG state, rollout position,
guardrail calibration, and health marker.\\
Final policy suite &
\path{rlhf_runs/qwen25_05b_helpsteer3_trl_a100_full/full/eval_ppo_guarded_r768_b64_kl10_rm_ep2_u188_full2017/}.
Authoritative 2,017-prompt summary, full rows, plots, and audit queues.\\
Full response artifact &
\path{rlhf_runs/portfolio_full_policy_comparisons_final_trl.json}. Redacted
Base/PPO explorer payload.\\
\bottomrule
\end{longtable}
\normalsize

The active command sequence is:

\begin{verbatim}
python -m scripts.rlhf_trl_prepare_data \
  --config configs/trl/qwen25_05b_helpsteer3_sft.yaml
python -m scripts.rlhf_trl_train_sft \
  --config configs/trl/qwen25_05b_helpsteer3_sft.yaml
python -m scripts.rlhf_trl_train_reward_model \
  --config configs/trl/qwen25_05b_helpsteer3_reward.yaml
python -m scripts.rlhf_trl_train_ppo \
  --config configs/trl/qwen25_05b_helpsteer3_ppo.yaml
python -m scripts.rlhf_evaluate_policy_suite \
  --config configs/trl/qwen25_05b_helpsteer3_eval_suite.yaml
python -m scripts.rlhf_audit_policy_suite \
  --eval-dir \
  rlhf_runs/qwen25_05b_helpsteer3_trl_a100_full/full/\
  eval_ppo_guarded_r768_b64_kl10_rm_ep2_u188_full2017 \
  --base-label base --sft-label sft_trl --ppo-label ppo_trl
python -m scripts.rlhf_sync_runs --mode lightweight
\end{verbatim}

The DPO preparation, doctor, training, and evaluation commands are retained in
\path{docs/REPOSITORY_GUIDE.md} and the DPO subsection of this companion. They
are not part of the executed final result. The executed guarded PPO notebook
and persisted resolved configuration are authoritative when a historical path
name conflicts with a runtime override.

\section{Metric Interpretation}

\begin{description}[style=nextline]
  \item[SFT mean token accuracy]
  Fraction of preferred completion tokens predicted correctly under teacher
  forcing; not response-level task accuracy.

  \item[Reward-model pairwise accuracy]
  Fraction of held-out pairs for which the chosen response scores above the
  rejected response.

  \item[Policy win rate]
  Fraction of evaluation prompts for which one policy receives the higher
  learned reward. It is a proxy-judge result, not a human preference rate.

  \item[Objective KL]
  Sampled response-level departure from the frozen SFT reference. It is not the
  same statistic as PPO's old-to-new update approximate KL.

  \item[Cap hit]
  The response uses the full evaluation generation budget and does not stop
  naturally before 768 new tokens.

  \item[Repeated 4-gram fraction]
  Fraction of word-level 4-gram positions attributable to repeated 4-grams,
  used only to prioritize qualitative review.
\end{description}

\begin{thebibliography}{99}

\bibitem{christiano2017}
P. F. Christiano, J. Leike, T. Brown, M. Martic, S. Legg, and D. Amodei.
\newblock Deep reinforcement learning from human preferences.
\newblock In \emph{Advances in Neural Information Processing Systems}, 2017.
\newblock \url{https://arxiv.org/abs/1706.03741}.

\bibitem{schulman2015}
J. Schulman, S. Levine, P. Abbeel, M. Jordan, and P. Moritz.
\newblock Trust Region Policy Optimization.
\newblock In \emph{Proceedings of the 32nd International Conference on Machine
Learning}, 2015.
\newblock \url{https://arxiv.org/abs/1502.05477}.

\bibitem{schulman2017}
J. Schulman, F. Wolski, P. Dhariwal, A. Radford, and O. Klimov.
\newblock Proximal Policy Optimization Algorithms.
\newblock arXiv:1707.06347, 2017.
\newblock \url{https://arxiv.org/abs/1707.06347}.

\bibitem{stiennon2020}
N. Stiennon, L. Ouyang, J. Wu, D. Ziegler, R. Lowe, C. Voss,
A. Radford, D. Amodei, and P. F. Christiano.
\newblock Learning to summarize with human feedback.
\newblock In \emph{Advances in Neural Information Processing Systems}, 2020.
\newblock \url{https://arxiv.org/abs/2009.01325}.

\bibitem{ouyang2022}
L. Ouyang et al.
\newblock Training language models to follow instructions with human feedback.
\newblock In \emph{Advances in Neural Information Processing Systems}, 2022.
\newblock \url{https://arxiv.org/abs/2203.02155}.

\bibitem{huang2024}
S. Huang, M. Noukhovitch, A. Hosseini, K. Rasul, W. Wang, and L. Tunstall.
\newblock The N+ Implementation Details of RLHF with PPO: A Case Study on
TL;DR Summarization.
\newblock arXiv:2403.17031, 2024.
\newblock \url{https://arxiv.org/abs/2403.17031}.

\bibitem{wang2025}
Z. Wang, J. Zeng, O. Delalleau, H.-C. Shin, F. Soares, A. Bukharin,
E. Evans, Y. Dong, and O. Kuchaiev.
\newblock HelpSteer3-Preference: Open Human-Annotated Preference Data across
Diverse Tasks and Languages.
\newblock arXiv:2505.11475, 2025.
\newblock \url{https://arxiv.org/abs/2505.11475}.

\bibitem{qwen25}
Qwen Team.
\newblock Qwen2.5 Technical Report.
\newblock arXiv:2412.15115, 2024.
\newblock \url{https://arxiv.org/abs/2412.15115}.

\bibitem{hu2022}
E. J. Hu et al.
\newblock LoRA: Low-Rank Adaptation of Large Language Models.
\newblock In \emph{International Conference on Learning Representations},
2022.
\newblock \url{https://arxiv.org/abs/2106.09685}.

\bibitem{trl}
L. von Werra et al.
\newblock TRL: Transformer Reinforcement Learning.
\newblock Hugging Face, 2020--2026.
\newblock \url{https://github.com/huggingface/trl};
\url{https://huggingface.co/docs/trl/ppo_trainer}.

\bibitem{gao2023}
L. Gao, J. Schulman, and J. Hilton.
\newblock Scaling Laws for Reward Model Overoptimization.
\newblock In \emph{Proceedings of the 40th International Conference on Machine
Learning}, 2023.
\newblock \url{https://arxiv.org/abs/2210.10760}.

\bibitem{rafailov2023}
R. Rafailov, A. Sharma, E. Mitchell, C. D. Manning, S. Ermon, and C. Finn.
\newblock Direct Preference Optimization: Your Language Model is Secretly a
Reward Model.
\newblock In \emph{Advances in Neural Information Processing Systems}, 2023.
\newblock \url{https://arxiv.org/abs/2305.18290}.

\end{thebibliography}

\end{document}
